\bookStart{Instigation of Guthrun}[Guð·rúnar hvǫt]
\setBookCode{Gudrunarhvot}

\begin{flushright}%
\textbf{Dating} \parencite{Sapp2022}: early C11th (0.781)--late C11th (0.177)

\textbf{Meter:} \Fornyrdislag
\end{flushright}%

\section{Introduction}

\subsection{Preservation}

The \textbf{Instigation of Guthrun} (\Gudrunarhvot) is the penultimate poem of \Regius, where it is found on fol. 44r/13--44v/27.  The poem is complete.

\subsection{Story}

As its name spells out, \Gudrunarhvot\ deals with Guthrun’s instigation of her two sons Sarle and Hamthew to avenge their sister Swanhild, killed by the cruel tyrant Ermenric.  After the sons depart on their inevitably fatal course, Guthrun laments her cruel fate before addressing her long-dead first husband and true love, Siward.  She constructs a great funeral pyre in order to join him, it is implied, by self-immolation.

A variant of Guthrun’s instigation and the eager response of her sons (\textlink{Gudrunarhvot}[2]--8), but not her following monologue, is also found in \textlink{Hamdismal}[3]--10.  These sections share so much in common, including numerous identical or near-identical lines (\textlink{Gudrunarhvot}[2]/3a, 4b--6 = \textlink{Hamdismal}[3]/2--4,
\textlink{Gudrunarhvot}[4]/1--3 = \textlink{Hamdismal}[6]/1--3,
\textlink{Gudrunarhvot}[4]/4 = \textlink{Hamdismal}[7]/1,
\textlink{Gudrunarhvot}[4]/5a \char`~= \textlink{Hamdismal}[7]/2b), that they must be related.  It is however clear from the variant forms that this relation is through the oral tradition (cf. e.g. \textlink{Gudrunarhvot}[4]/3 and \textlink{Hamdismal}[6]/3, which are near-identical, but where the former appears to be an unmetrical corruption of a line originally resembling the latter; in a scribal context such corruption could never happen).  Since both poems are related through oral tradition, this of course means that neither poem can be a scribal composition.

\subsubsection{Historical basis}

Because both \Gudrunarhvot\ and the immediately following \textlink{Hamdismal} deal with the same events it is best to discuss their historical basis in a single section.

It is a remarkable thing that the narrative described in these two poems is surprisingly close to historical fact, or at least what we know as such.  In Jordanes’ \emph{Getica} 129 we read the following:

\begin{quote}\begin{small}
  \emph{Nam Hermanaricus, rex Gothorum, licet, ut superius retulimus, multarum gentium extiterat triumphator, de Hunnorum tamen adventu dum cogitat, Rosomonorum gens infida, quae tunc inter alias illi famulatum exhibebat, tali eum nanciscitur occasione decipere. Dum enim quandam mulierem Sunilda nomine ex gente memorata pro mariti fraudulento discessu rex furore commotus equis ferocibus inligatam incitatisque cursibus per diversa divelli praecipisset, fratres eius Sarus et Ammius, germanae obitum vindicantes, Hermanarici latus ferro petierunt; quo vulnere saucius egram vitam corporis inbecillitate contraxit.}

  ‘Now although Hermanaric, king of the Goths, was the conqueror of many tribes, as we have said above, yet while he was deliberating on this invasion of the Huns, the treacherous tribe of the Rosomoni, who at that time were among those who owed him their homage, took this chance to catch him unawares. For when the king had given orders that a certain woman of the tribe I have mentioned, Sunilda by name, should be bound to wild horses and torn apart by driving them at full speed in opposite directions (for he was roused to fury by her husband’s treachery to him), her brothers Sarus and Ammius came to avenge their sister’s death and plunged a sword into Hermanaric’s side. Enfeebled by this blow, he dragged out a miserable existence in bodily weakness.’ \parencite[87]{Jordanes1915}
\end{small}\end{quote}

It is known that Ermenric died in 375 \ce\ and that Jordanes was working ca. 550 \ce.  So it is that in two poems composed at least 450 years (probably closer to 550) after Ermenric’s death, preserved in a manuscript written 900 years after his death, we find the very same events retold---the torture-murder of Swanhild with horses, the execution of her lover, the revenge of her two brothers Sarle and Hamthew, their maiming of the king---events retold not on the basis of historical learning, but rather through an independent oral tradition.

\subsection{Style}

The poem is rather sparse with kenning and enjambment.  Many sts.~are very touching, e.g. 13, 15--16, and 19.

\subsection{Meter}

The meter is regular \Fornyrdislag.

\sectionline

\newpage

\section{Text}

\subsection{Frá Guð·rúnu (‘From Guthrun’)}

\bpg\bpa\mssnote{\Regius~44r/13--20}%
\edtext{Guð·rún}{\Afootnote{The \emph{G} is an ornamented capital 2 lines high in \Regius.}} gekk þá til sę́var er hón hafði drepit Atla, gekk út á sę́’inn ok vildi fara sér. Hón mátti eigi søkkva. Rak hana yfir fjǫrð’inn á land Jón·akrs konungs; hann fekk hennar. Þeira synir vóru þeir Sǫrli ok Erpr ok Ham·ðir. Þar fǿddisk upp Svan·hildr Sig·urðar dóttir. Hón var gift Jǫrmun·rekk inum ríkja. Með hánum var Bikki. Hann réð þat at Rand·vér konungs son skyldi taka hana; þat sagði Bikki konungi. Konungr lét hengja Rand·vé en troða Svan·hildi undir hrossa fótum. En er þat spurði Guð·rún þá kvaddi hón sonu sína.\epa

\bpb {\huge G}\textsc{uthrun then went to the sea} after she had slain Attle, went out into the sea and would end herself.  She could not sink.  She was driven across the firth to the land of king Enacker; he got her for wife.  Their sons were Sarle and Earp and Hamthew.  There Swanhild, Siward’s daughter, was reared.  She was married off to Ermenric the powerful.  With him was \inx[P]{Bicke}.  He counseled that Randwigh, the king’s son, should sleep with her; this Bicke told the king.  The king let Randwigh be hanged, but Swanhild he had trampled beneath the hooves of horses.  But when Guthrun learned of that she called out to her sons.\epb\epg

\sectionline

\subsection{Guð·rúnar hvǫt (‘The Instigation of Guthrun’)}

\bvg\bva%
\edtext{Þá}{\Afootnote{The \emph{Þ} is an ornamented initial 3 lines high, with the left vertical stroke ascending 3 lines and descending 2 lines in the margin in \Regius.}} frá’k \alst{s}ęnnu \hld\ \alst{s}líðr-fęng⸗ligasta, &
\alst{t}rauð mǫ́l \alst{t}alit \hld\ af \alst{t}rega stórum, &
es \alst{h}arð-\alst{h}uguð \hld\ \alst{h}vatti at vígi &
\alst{g}rimmum orðum \hld\ \alst{G}uð·ru̇n sonu:\eva

\bvb {\huge T}\textsc{hat gibing I’ve found} most direly caught--- \\
hard-pressed speeches told from great grief--- \\
when hard-hearted she goaded to war, \\
with fierce words, Guthrun, her sons:\evb\evg


\bvg\bva%
„Hví \alst{s}itið \edtext{\emph{it}}{\Afootnote{metr. emend.; om. by haplography (\emph{*‘ſıtıt ıt’} > \emph{‘ſıtıt’}) \Regius}}? \hld\ Hví \alst{s}ofið lífi? &
Hví \alst{t}regr⸗at ykkr \hld\ \alst{t}ęiti at mę́la, &
\edtext{es \alst{Jǫ}rmun·rekr \hld\ \alst{y}ðra systur, &
\alst{u}nga at \alst{a}ldri, \hld\ \alst{jó}m of traddi, &
\alst{h}vítum ok svǫrtum \hld\ ȧ \alst{h}ęr-vegi &
\alst{g}rǫ́m, \alst{g}ang-tǫmum \hld\ \alst{G}otna hrossum?}{\lemma{es \dots\ hrossum. ‘when \dots\ horses?’}\Bfootnote{Repeated almost identically in \textlink{Hamdismal}[3].}}\eva

\bvb “Why do ye two sit? Why sleep ye your lives away? \\
Why does it not trouble you to speak merrily \\
when Ermenric has had your sister \\
young of age trampled by steeds, \\
by the white and black on the path of war, \\
by the grey, pacing, Gotnish horses?\evb\evg


\bvg\bva%
Urðu⸗a it \alst{g}·líkir \hld\ þęim \alst{G}unn·ari &
né in \alst{h}ęldr \alst{h}ugðir \hld\ sęm vas \alst{H}ǫgni; &
\alst{h}ęnnar mynduð it \hld\ \alst{h}ęfna lęita &
ef it \alst{m}óð ę́ttið \hld\ \alst{m}ïnna brǿðra &
eða \alst{h}arðan \alst{h}ug \hld\ \alst{H}u̇n-konunga.“\eva

\bvb Ye have not turned out alike to Guther \\
nor instead thoughtful like Hain was. \\
\emph{Her} would ye have sought to avenge \\
if ye had the courage of my brothers \\
or the hard heart of the Hunnish kings!”\evb\evg


\bvg\bva%
\Ballnote{Closely corresponding to \textlink{Hamdismal}[6]--7/2; ll. 1--4 and 5b are found in related variant forms there.}%
Þȧ kvað þat \alst{H}am·ðir \hld\ inn \alst{h}ugum stóri: &
„\alst{L}ítt myndir þú \hld\ \alst{l}ęyfa dǫ́ð Hǫgna &
\edtext{þȧ’s \alst{S}ig·urð vǫkðu \hld\ \alst{s}vefni ór}{\Bfootnote{This l.~is metrically defective in the present poem but survives better as \textlink{Hamdismal}[6]/3.}}; &
\alst{b}ǿkr vǫ́ru þïnar, \hld\ inar \alst{b}lá-hvítu, &
roðnar ï \alst{v}ers dręyra, \hld\ folgnar ï \alst{v}al-blóði.\eva

\bvb Then this quoth Hamthew the great of heart: \\
“Thou hadst little cause to praise Hain’s deeds \\
when they awoke Siward from his sleep. \\
Thy blue-white bed-clothes were \\
reddened in thy husband’s gore, buried in death-blood.”\evb\evg


\bvg\bva%
Urðu þér [...] \hld\ brǿðra hęfndir &
\alst{s}líðrar ok \alst{s}árar \hld\ es þú \alst{s}onu myrðir; &
knę́ttim [...] \hld\ Jǫrmun·rekki, &
\alst{s}am-hyggjendr, \hld\ \alst{s}ystur hęfna.\eva

\bvb For thee thy [...] deeds of revenge for thy brothers became \\
biting and sore when thou murderedst thy sons. \\
We shall [...] with a single mind \\
avenge our sister on Ermenric.\evb\evg


\bvg\bva%
Berið \alst{h}nossir framm \hld\ \alst{h}u̇n-konunga, &
\alst{h}ęfir þú okkr \alst{h}vatta \hld\ at \alst{h}jǫr-þingi.“\eva

\bvb Bring forth the treasures of the Hunnish kings! \\
Thou hast incited us to the thing of swords \ken{battle}.”\evb\evg


\bvg\bva%
\alst{H}lę́jandi Guð·ru̇n \hld\ \alst{h}varf til skęmmu, &
\alst{k}umbl \alst{k}onunga \hld\ ór \alst{k}ęrum valði, &
\alst{s}íðar brynjur \hld\ ok \alst{s}onum fǿrði; &
hlóðu⸗sk \alst{m}óðgir \hld\ ȧ \alst{m}ara bógu.\eva

\bvb Laughing, Guthrun turned to her chamber. \\
The heirlooms of kings from the chests she picked, \\
the trailing byrnies, and brought them to her sons. \\
Fierce, they loaded themselves onto the backs of steeds.\evb\evg


\bvg\bva%
Þȧ kvað þat \alst{H}am·ðir \hld\ inn \alst{h}ugum stóri: &
„Svá \edtext{kom⸗a’\emph{k}}{\Afootnote{\emph{‘comaz’} \Regius}} \alst{m}ęirr aptr \hld\ \alst{m}óður at vitja &
\alst{g}ęir-Njǫrðr hniginn \hld\ ȧ \alst{G}oð-þjóðu &
at þú \alst{ę}rfi \hld\ at \alst{ǫ}ll oss drykkir, &
at \alst{S}van·hildi \hld\ ok \alst{s}onu þïna.“\eva

\bvb Then this quoth Hamthew the great of heart: \\
“So will I no more come back to visit my mother, \\
{[when I am]} a spear-Nearth \ken{warrior} fallen in the land of the Gots, \\
after thou hast drunk a wake for us all, \\
for Swanhild and for thy sons.”\evb\evg


\bvg\bva%
\Ballnote{The two boys depart (cf. \textlink{Hamdismal}), and Guthrun’s long monologue begins, lasting from st. 10 to 20.}%
\alst{G}uð·ru̇n \alst{g}rátandi, \hld\ \alst{G}júka dóttir, &
gekk \alst{t}reg⸗liga \hld\ ȧ \alst{t}ȧi sitja &
ok at \alst{t}ęlja, \hld\ \alst{t}ǫ́rug-hlýra, &
\alst{m}óðug spjǫll \hld\ ȧ \alst{m}argan veg:\eva

\bvb Guthrun, weeping, Yivick’s daughter, \\
went in grief to sit on the ground, \\
and to tell with teary cheeks \\
her moving tale in many ways:\evb\evg


\bvg\bva%
„Þrjȧ vissa’k \alst{ę}lda, \hld\ þrjȧ vissa’k \alst{a}rna, &
\alst{v}as’k þrimr \alst{v}erum \hld\ \alst{v}egin at húsi; &
\alst{ęi}nn vas mér Sig·urðr \hld\ \alst{ǫ}llum bętri &
es \alst{b}rǿðr mïnir \hld\ at \alst{b}ana urðu.\eva

\bvb “I have known three fires; I have known three hearths; \\
for three husbands I have been brought to the house. \\
To me Siward alone was better than all--- \\
he whose bane my brothers became.\evb\evg


\bvg\bva%
\edtrans{\alst{S}várra \alst{s}ára \hld\ \alst{s}á’k⸗at, né kunn\emph{a}}{Heavy griefs I neither saw nor knew}{\Bfootnote{Before the killing of Siward.  Guthrun thus presents her life as idyllic before her brothers killed her lover.}} &
\edtext{[...]}{\Bfootnote{No gap is indicated in \Regius, but it seems inevitable that a line has fallen out here containing something to the effect of “before my brothers murdered Siward”.}} &
\edtrans{\alst{m}ęirr}{moreover}{\Bfootnote{Namely after the killing of Siward, when she was married off to Attle against her will.}} þȯttu⸗sk \hld\ \alst{m}ér of stríða &
es mik \alst{ǫ}ðlingar \hld\ \alst{A}tla gǫ́fu.\eva

\bvb Heavy griefs I neither saw nor knew, \\
{[...]}; \\
moreover I thought it injurious to me \\
when the athlings gave me to Attle.\evb\evg


\bvg\bva%
\Ballnote{Guthrun alludes to her murder of her children with Attle, as dealt with in \textlink{Atlakvida}[37]--42 and \textlink{Atlamal}[73]--81.  Cf. \textlink{Hamdismal}[8].}%
\edtrans{\alst{H}u̇na \alst{h}vassa}{The smart bear-cubs}{\Bfootnote{Her sons with Attle; for the appellation cf. \textlink{Volundarkvida}[24]/1.}} \hld\ \alst{h}ét’k mér at ru̇num; &
mátti’g⸗a’k \alst{b}ǫlva \hld\ \alst{b}ǿtr of vinna &
áðr ek \alst{h}nóf \alst{h}ǫfuð \hld\ af \edtext{\emph{\alst{H}}niflungum}{\Afootnote{metr. emend. (cf. \textlink{HelgakvidaOne}[48]/5); \emph{Niflungum} \Regius}}.\eva

\bvb The smart bear-cubs I called to me for counsel; \\
I could not find remedies for the bales \\
before I severed the heads from the Nivlings.\evb\evg


\bvg\bva%
\Ballnote{Guthrun attempted to take her own life, but was instead driven to the realm of king Enacker, who married her.  See P1 above.}%
\alst{G}ekk ek til strandar, \hld\ \alst{g}rǫm vas’k nornum, &
vilda’k hrinda \hld\ stríð grið þęira; &
\edtrans{\alst{h}ófu mik, né drękkðu, \hld\ \alst{h}ǫ́var bǫ́rur}{I was lifted---not drowned---by the billows high}{\Bfootnote{Clearly related to the prophecy of this event in \textlink{Sigurdskamma}[62]/3, with which the b-verse is identical.}}, &
því \alst{l}and of sté’k \hld\ at \alst{l}ifa skylda’k.\eva

\bvb I went to the shore; I was wroth with the norns; \\
I wished to break their stubborn peace. \\
I was lifted---not drowned---by the billows high, \\
so I stepped aland since I was meant to live.\evb\evg


\bvg\bva%
Gekk ek ȧ \alst{b}ęð \hld\ ---hugða’k mér fyr \alst{b}ętra--- &
\alst{þ}riðja sinni \hld\ \alst{þ}jóð-konungi; &
\alst{ó}l ek mér \alst{jó}ð, \hld\ \alst{ę}rfi-vǫrðu &
{[...]} \hld\ Jón·akrs \edtext{sonu}{\Bfootnote{emend.; \emph{sonum} \Regius}}.\eva

\bvb I mounted the bed---I thought it better for me--- \\
of a great king for a third time. \\
I begot babes, heritance-defenders \ken{sons}, \\
{[...]}, Enacker’s sons.\evb\evg


\bvg\bva%
Ęn umb \alst{S}van·hildi \hld\ \alst{s}ǫ́tu þýjar, &
es mïnna \alst{b}arna \hld\ \alst{b}atst full-hugða’k; &
\alst{s}vá vas \alst{S}van·hildr \hld\ ï \alst{s}al mïnum, &
\alst{s}ęm vę́ri \alst{s}ǿm-lęitr \hld\ \alst{s}ólar gęisli.\eva

\bvb But around Swanhild handmaids sat, \\
she whom I thought best of among my children. \\
So was Swanhild in my hall \\
like were she a beautiful ray of the sun.\evb\evg


\bvg\bva%
\alst{G}ǿdda’k \alst{g}olli \hld\ ok \alst{g}uð-vęfjum, &
áðr ek \alst{g}ę́fa’k \hld\ \alst{G}oð-þjóðar til; &
þat ’s mér \alst{h}arðast \hld\ \alst{h}arma mïnna &
of þann inn \alst{h}víta \hld\ \alst{h}add Svan·hildar: &
\alst{au}ri trǫddu \hld\ und \alst{jó}a fótum;\eva

\bvb I endowed her with gold and with god-weave \\
before I gave her to the folk of the Gots. \\
It is for me the hardest of my harms, \\
over the white hair of Swanhild--- \\
they trampled it in mud beneath the hooves of steeds!---\evb\evg


\bvg\bva%
ęn sá \alst{s}árastr, \hld\ es þęir \alst{S}ig·urð mïnn, &
\alst{s}igri rę̇ntan, \hld\ ï \alst{s}ę́ing vǫ́gu; &
\edtext{ęn sá \alst{g}rimmastr, \hld\ es þęir \alst{G}unn·ari, &
\alst{f}rȧnir ormar, \hld\ til \alst{f}jǫrs skriðu; &
ęn sá \alst{h}vassastr, \hld\ es til \alst{h}jarta &
\alst{k}onung ȯ·blauðan \hld\ \alst{k}vikvan skǫ́ru.}{\lemma{ęn \dots\ skǫ́ru ‘but \dots\ alive.’}\Bfootnote{For the deaths of Guther and Hain cf. \textlink{Atlakvida}[23]--35, \textlink{Atlamal}[56]--64.}}\eva

\bvb but the sorest when my Siward, \\
robbed of victory they slew in bed, \\
but the cruelest when for Guther \\
the gleaming serpents slithered to his bane, \\
but the sharpest when unto the heart \\
they cut the unsoft king \ken*{= Hain}, alive.\evb\evg


\bvg\bva%
\Ballnote{After recounting her cruel fates, Guthrun addresses Siward, who is long dead.  The tone is at last one of release.  With the departure of her doomed sons Guthrun has nothing to live for upon the earth; her family line is ended, as Byrnhild predicted (\textlink{Sigurdskamma}[63]--64) and she now knows that at last her hard life is over.  She will at last find happiness again, but only in death, reunited with her true love Siward.---
An interesting question not dealt with in the poetry is the implication of the fact that Byrnhild has already joined Siward by suttee immediately after his death (\textlink{Sigurdskamma}[65]--71, \textlink{Helreid}).  Do the two rival women encounter each other in the afterlife as co-wifes, or does one oust the other?}%
Fjǫlð man’k bǫlva, \hld\ \edtext{[...]}{\Bfootnote{A half-line is undoubtedly missing here.  It would presumably have mirrored the first half-line, containing \emph{fjǫlð man’k b-} ‘I recall a multitude of ...’ with alliteration on a noun beginning with \emph{b}.}} &
\alst{b}ęit-tu, Sig·urðr, \hld\ inn \alst{b}lakka mar, &
\alst{h}ęst inn \alst{h}rað-fǿra \hld\ lát-tu \alst{h}inig renna! &
\alst{S}itr ęigi hér \hld\ \alst{s}nǫr né dóttir &
sú’s \alst{G}uð·ru̇nu \hld\ \alst{g}ę́fi hnossir.\eva

\bvb I recall a multitude of bales; [...] \\
Saddle, Siward, the fallow steed, \\
the quick-pacing horse---let him run hither! \\
Here sits no son’s wife nor daughter \\
who to Guthrun might give treasures.\evb\evg


\bvg\bva%
\alst{M}inns-tu, Sig·urðr, \hld\ hvat vit \alst{m}ę́ltum &
þȧ’s vit ȧ \alst{b}ęð \hld\ \alst{b}ę́ði sǫ́tum? &
at þú \alst{m}yndir \alst{m}ïn \hld\ \alst{m}óðugr vitja, &
\alst{h}alr, ór \alst{h}ęlju, \hld\ en ek þïn ór \alst{h}ęimi.\eva

\bvb Dost thou remember, Siward, what we said \\
when on the bed we both did sit, \\
that thou, fierce man, wouldst visit me \\
from Hell---and I thee from the world?\evb\evg


\bvg\bva%
\Ballnote{Guthrun orders the construction of her own pyre.  It is not necessary to think that she is burned alive, although that is likely.}%
Hlaðið \alst{é}r, \alst{ja}rlar, \hld\ \alst{ęi}ki-kǫst’inn, &
látið þann und \edtrans{\alst{h}i\emph{mn}i}{heaven}{\Afootnote{emend.; \emph{hilmi} ‘prince’ \Regius}\Bfootnote{Although \emph{und hilmi} could be read as ‘around the prince’, the pyre is certainly Guthrun’s.  At this point in the narrative Siward is long dead, and the point of the pyre is for Guthrun to join him in the afterlife.  Since no prince is being burned this requires emendation, and \emph{himni} ‘heaven’ (dat.~sg.) easily presents itself.  It is a common poetic motif for something to be ‘the best under heaven’; cf. \textlink{Gripisspa}[10]/4, \Ynglingatal\ 27, etc. (TODO.)}} \hld\ \alst{h}ę́stan verða! &
Męgi \alst{b}ręnna \edtext{\alst{b}rjóst \hld\ \alst{b}ǫlva-fullt}{\lemma{brjóst bǫlva-fullt ‘my curse-filled breast’}\Bfootnote{The possessive pronoun is implicit.  With these words Guthrun reflects on the many evils she has been responsible for (cf. \textlink{Atlakvida}[47]).}} ęldr &
umb hjarta [\dots] \hld\ þiðni sorgir!“\eva

\bvb Load, ye earls, the oaken pile \ken{pyre}! \\
Let it become the highest under heaven! \\
May fire burn my curse-filled breast, \\
unto the heart \dots\ may the sorrows melt away!”\evb\evg


\bvg\bva%
\Ballnote{The last st. steps out of the world of the poem and should be read in the voice of the narrating ballad-poet, who indirectly addresses the members of the audience and wishes them better fortunes than the character in the poem.  The aristocratic titles afforded to the audience are notable.}%
\alst{Jǫ}rlum \alst{ǫ}llum \hld\ \alst{ó}ðal batni, &
\alst{s}nótum ǫllum \hld\ \alst{s}org at minni &
at þetta \alst{t}reg-róf \hld\ of \alst{t}alit vę́ri.\eva

\bvb For all earls may patrimony improve--- \\
for all ladies may sorrow decrease \\
after this grief-chain has been recounted!\evb\evg

\sectionline
